\subsection{Twisted cubic Cartier blocks and the theta embedding}
\label{subsec:ppav-twisted-Cartier-Moore}

Work over an algebraically closed field of characteristic two.  All
Cartier maps are relative.  This subsection uses the complete principal
addition proof and the separately certified universal Moore determinant
classification; neither a characteristic-zero theta identity nor an
untyped identification of Frobenius twists is used.

\begin{lemma}[The principal coefficient on an arbitrary abelian surface]
\label{lem:principal-theta-Cartier-translated-square}
\status{Proved; independent audits passed.}
Let \(B\) be any abelian surface and let \(M\) be symmetric and
principally polarizing.  Fix \(0\ne\theta\in H^0(B,M)\).
There exist \(b_0\in B[2](k)\), an isomorphism
\(t_{b_0}^*M^2\simeq M^2\), and a nonzero scalar \(c\), such that
\[
 C_z\bigl(\theta(z+a)\theta(z-a)\omega_z\bigr)
   =c\,\theta'(z)\,\mathcal T_{b_0}(\theta^2)(a)\,\omega_z'.
                                                        \tag{TCM1}
\]
The product is interpreted through the symmetric theorem-of-square
isomorphism.  Here \(\mathcal T_{b_0}\) is the translation operator
on \(M^2\) defined by the stated isomorphism.  In two-rank zero
one has \(b_0=0\), recovering the preceding principal formula.
\end{lemma}
\begin{proof}
The addition tensor is nondegenerate by
Lemma~\ref{lem:principal-theta-addition-perfect}, whose proof uses no
two-rank assumption.  The global principal Cartier trace is also
surjective without that assumption.  Indeed, put \(G=\ker(F_B)\).
The dual algebra \(\mathcal O(G)^*\) is a commutative Artinian
algebra, so every simple representation is one-dimensional, given
by a character of \(G\).  The trace kernel, after pullback to the
proper torsor \(B\), is a constant subrepresentation of the regular
representation, exactly as in
Lemma~\ref{lem:principal-theta-Cartier-surjective}.  Its composition
series descends to line bundles \(Q\) on \(B^{(2)}\) with
\(F_B^*Q\simeq\mathcal O_B\).  They are two-torsion, hence
numerically trivial.  Tensoring by \(M^{(2)}\) makes all factors
ample, with zero higher cohomology.  The same exact-sequence argument
therefore proves surjectivity of the global trace.

Write the Cartier contraction of the addition tensor as
\(\theta'(z)\kappa(a)\omega_z'\), with
\(0\ne\kappa\in H^0(B,M^2)\).  Nonvanishing follows from
nondegeneracy and the just-proved trace surjectivity.  The natural
and canonical diagonal \(G\)-linearizations differ by a character
\(\chi\), by the proper-base argument in
Proposition~\ref{prop:principal-theta-Cartier-square}.
Trace equivariance implies that \(\kappa\) has the corresponding
character for the canonical \(G\)-action.  Descent thus identifies
its one-dimensional eigenspace with
\(H^0(B^{(2)},M^{(2)}\otimes Q_\chi)\), for a line bundle
\(Q_\chi\in\ker(F_B^*)\).  The inverse-character convention is
immaterial here, since these characters and line bundles have order
at most two.

The principal polarization of \(M^{(2)}\) gives a unique
\(b_0'\in B^{(2)}[2](k)\) with
\(M^{(2)}\otimes Q_\chi\simeq t_{b_0'}^*M^{(2)}\).
Relative Frobenius is a universal homeomorphism and is bijective
on geometric points; let \(b_0\) be the unique geometric lift of
\(b_0'\).  Since \(F_B(2b_0)=0\) and \(\ker F_B\) is connected,
\(2b_0=0\).  Pulling the descended theta section back gives precisely
a nonzero scalar times \(\mathcal T_{b_0}(\theta^2)\), proving
(TCM1).  If the two-rank is zero, there is only one geometric
two-torsion point, so \(b_0=0\).
\end{proof}

\begin{lemma}[A symmetric level-three addition normalization]
\label{lem:ppav-symmetric-level3-addition-gauge}
\status{Proved; independent audits passed.}
Let \((A,\Theta_A)\) be a principally polarized abelian surface,
with a symmetric representative, and put \(L=\mathcal O_A(3\Theta_A)\).
Choose complementary Lagrangians \(K,K'\subset A[3]\), with
\(K'\simeq H:=\mathbf F_3^2\).  One may choose order-three theta
splittings over both groups compatible with normalized symmetry.
The quotient \(\pi:A\to B=A/K\) then has a symmetric principal
line bundle \(M\) with \(\pi^*M\simeq L\).  If
\(e_0=\pi^*\theta\) and \(e_i=R_i e_0\) for the chosen
translation operators \(R_i\), then \((e_i)_{i\in H}\) is a theta
basis and normalized symmetry sends \(e_i\) to \(e_{-i}\).

For the rigidified symmetric square isomorphism on \(A\times A\),
the lift of translation \((k,h)\in K'\times K'\) obtained by
pulling back \(R_{k+h}\boxtimes R_{k-h}\) agrees exactly with
\(R_k^2\boxtimes R_h^2\).
\end{lemma}
\begin{proof}
The prime-to-characteristic theta-group structure and maximal
isotropic descent give splittings and a principal descendant
\citep[Corollary~8.20, Lemma~8.24, Theorem~8.28;
Theorem~8.10 and Corollary~8.14]{EdixhovenVanDerGeerMoonenAV}.
Here and below a theta translation operator means an isomorphism
\(t_t^*L\to L\), acting on sections by translation followed by that
isomorphism; its central scalar acts with weight one.

Let \(\sigma\) denote the involution on the theta group induced by
the normalized symmetry of \(L\).  For any splitting \(R_t\) of
one of the odd-order Lagrangians, write
\(\sigma(R_t)=\eta(t)R_{-t}\); its discrepancy \(\eta\) is a
character.  Replacing \(R_t\) by \(\psi(t)R_t\), with
\(\psi^2=\eta^{-1}\), makes it symmetric.  Such \(\psi\) exists
uniquely because the character group has odd order.  The symmetric
splitting on \(K\) descends normalized symmetry to \(M\).
The unique theta section on \(B\) is fixed by this symmetry: its
scalar squares to one, hence equals one in characteristic two.
Thus \(\sigma e_i=e_{-i}\).  The commutator pairing identifies the
nine distinct \(K\)-characters of the \(e_i\), so they form a basis.
The normalized symmetry and its square being the identity can also
be read directly from its normalization at zero
\citep[Definition~11.44, p.~178]{EdixhovenVanDerGeerMoonenAV}; no
unwritten proof in Lemma~11.45 is being imported.

It remains to verify the claimed exact scalar normalization.
Use the symmetric square isomorphism rigidified at \((0,0)\);
the theorem of the cube supplies it
\citep[Theorem~2.7 and Corollary~2.8, pp.~19--20]
{EdixhovenVanDerGeerMoonenAV}.
The two \(K'\times K'\)-linearizations in the statement have the
same base action.  Their ratio is a character \(\delta(k,h)\),
because every global unit on the proper connected product is a
scalar.  The rigidified square isomorphism is compatible with
each sign change \((z,a)\mapsto(-z,a)\) and
\((z,a)\mapsto(z,-a)\).  Indeed, on the target of
\((z,a)\mapsto(z+a,z-a)\) these are respectively interchange
followed by both normalized symmetries, and interchange alone;
the resulting isomorphisms agree at \((0,0)\), so agree globally.
The symmetry-compatible splittings consequently give
\[
 \delta(-k,h)=\delta(k,h),\qquad
 \delta(k,-h)=\delta(k,h).
\]
Writing the character as a product of its two restrictions shows
that both restrictions have order dividing two.  They also have
order dividing three, so \(\delta=1\).  This proves equality of
the actual lifts, not merely equality of their base translations.
\end{proof}

\begin{proposition}[The actual cross block is the Moore matrix]
\label{prop:ppav-Cartier-cross-Moore-identification}
\status{Proved; independent audits passed.}
Retain the preceding data.  There are a fixed \(a_0\in A[2](k)\)
and a nonzero scalar \(c\) such that the following holds.  For
\(a\in A(k)\), put \(P=\phi_L(a)\), so
\(t_a^*L\simeq L\otimes P\) and
\(t_{-a}^*L\simeq L\otimes P^{-1}\).
After choosing one fibre trivialization of the square isomorphism,
the nine-by-nine matrix of
\[
 H^0(A,LP)\times H^0(A,LP^{-1})
 \longrightarrow H^0(A^{(2)},L^{(2)}),\quad
 (s,t)\longmapsto C_{\rm rel}(st\omega_A)/\omega_A'
                                                        \tag{TCM2}
\]
evaluated by a variable functional is, up to one common nonzero
scalar and permutations of indices,
\[
 M_{c(a)}(x)_{u,v}=c(a)_{u-v}x_{u+v},\qquad
 [c(a)_h]_{h\in H}=[e_h(a+a_0)^2]_{h\in H}.             \tag{TCM3}
\]
In particular its parameter vector is the coordinatewise Frobenius
of the \(3\Theta_A\) embedding, evaluated at \(a+a_0\).
\end{proposition}
\begin{proof}
Choose \(\omega_A=\pi^*\omega_B\), and likewise on the Frobenius
twists.  The formula (TCM1) on \(B\), pulled back by the etale
isogeny in both variables, proves
\[
 C_z\bigl(e_0(z+a)e_0(z-a)\omega_z\bigr)
    =c\,e_0'(z)\,\mathcal T_{a_0}(e_0^2)(a)\,\omega_z'.
                                                        \tag{TCM4}
\]
Here \(a_0\) is the unique geometric two-torsion lift of \(b_0\)
through \(\pi\): an odd isogeny is an isomorphism on the full
two-torsion group schemes.  The operator \(\mathcal T_{a_0}\)
on \(L^2\) is the pullback of the operator on \(M^2\).

Apply the exactly normalized translation \((k,h)\) from the
preceding lemma to (TCM4).  Its input is
\(e_{k+h}(z+a)e_{k-h}(z-a)\).  On the Cartier target the squared
lift \(R_k^2\) descends to the Frobenius-twisted lift \(R_k^{(2)}\),
so its output is \(e_k'(z)\).  In the second variable
\(R_h^2\) commutes with \(\mathcal T_{a_0}\): their commutator
is the theta pairing on the two-torsion point \(a_0\) and the
three-torsion point \(h\), hence has order dividing both two and
three and equals one.  The second factor is therefore
\(\mathcal T_{a_0}(e_h^2)(a)\).

Consequently the fully typed global identity is
\[
 C_z\bigl(e_i(z+a)e_j(z-a)\omega_z\bigr)
   =c\,e_{(i+j)/2}'(z)\,
       \mathcal T_{a_0}(e_{(i-j)/2}^2)(a)\,\omega_z'.     \tag{TCM5}
\]
Division by two here is in \(H=\mathbf F_3^2\).
At a fixed \(a\), the operator \(\mathcal T_{a_0}\) multiplies
all evaluated squares by the same nonzero fibre isomorphism
\((L^2)_{a+a_0}\to(L^2)_a\).  A choice of trivialization gives
the projective vector asserted in (TCM3).
Finally reindex \(i=2u\), \(j=2v\), and let \(x_w\) be the
functional value on \(e_w'\).  Formula (TCM5) becomes exactly
(TCM3).  Input translations identify the two twisted section
spaces, and the square isomorphism identifies their product with
\(L^2\) after a single scalar fibre choice; no unrelated row-wise
or parameter-dependent coordinate identifications are assumed.
\end{proof}

\begin{theorem}[Perfect cross blocks away from two-torsion]
\label{thm:ppav-non-two-torsion-cross-Cartier}
\status{Proved; independent audits passed.}
For every principally polarized abelian surface in characteristic
two and every \(P\in\operatorname{Pic}^0(A)(k)\) with
\(P^2\not\simeq\mathcal O_A\), the cross pairing (TCM2) admits
a functional making its nine-by-nine matrix invertible.
\end{theorem}
\begin{proof}
Choose \(a\) with \(\phi_L(a)=P\), which is possible because
\(\phi_L\) is an isogeny.  By
Theorem~\ref{thm:DE-Moore-singular-parameters}, an identically singular
matrix (TCM3) would put its parameter vector in one of the eighty-one
translated character-scaled alternating three-spaces.
This set of spaces is preserved by coordinatewise Frobenius: the
primitive cube root is merely squared, permuting the characters.
Thus the unsquared evaluation vector \([e_h(a+a_0)]\) lies in
one of the same spaces.

The base space has equations \(c_0=0\), \(c_h=c_{-h}\), and is
contained in the projective fixed locus of the index involution
\(h\mapsto-h\).  The other eighty spaces are its translates by
the finite Heisenberg group: translation of the missing center and
character multiplication give all of them, since multiplication
by two is invertible on the character group.  In the symmetric
theta basis the base involution is the action of \([-1]_A\), and
the Heisenberg projective actions are exactly translations by
\(A[3]\).  Each space is therefore contained in the projective
fixed locus of the action induced by
\(x\mapsto2t-x\), for some \(t\in A[3](k)\).

The complete linear system \(|3\Theta_A|\) is a closed immersion.
For an indecomposable polarization this follows from
Theorem~\ref{thm:ppav-three-theta-jets}; for a product polarization,
the degree-three embeddings of the two elliptic factors followed
by the Segre embedding give the same assertion.  A fixed
evaluation point must consequently satisfy
\(2(a+a_0)=2t\).  Since \(2a_0=0\), this gives \(2a=2t\).
Applying \(\phi_L\), which kills \(A[3]\), yields
\(2P=0\), contrary to the hypothesis.  Hence the determinant
polynomial is nonzero, and some functional makes the cross block
invertible.  Notice that membership in the three-spaces is used
only to imply this fixed-point condition; equality with the whole
fixed-point locus is neither needed nor asserted.
\end{proof}

\begin{corollary}[Odd etale pullbacks of cubic principal polarizations]
\label{cor:ppav-odd-cover-maximal-dual-Cartier}
\status{Proved; independent audits passed.}
Let \(B\) be a nonordinary principally polarized abelian surface
in characteristic two, let \(M=\mathcal O_B(3\Theta_B)\), and let
\(\pi:A\to B\) be any odd-degree etale isogeny, of degree \(n\).
Put \(L=\pi^*M\).  The family of alternating forms
\[
 B_\lambda(s,t)=\lambda\bigl(C_{\rm rel}(st\omega_A)\bigr),
 \qquad \lambda\in H^0(A^{(2)},L^{(2)})^*,                 \tag{TCM6}
\]
has generic rank \(9n-1\) on \(H^0(A,L)\).
\end{corollary}
\begin{proof}
The odd kernel is constant over \(k\), and its character
decomposition gives
\(\pi_*\mathcal O_A=\bigoplus_\chi P_\chi\).
Every summand \(H^0(B,MP_\chi)\) has dimension nine.  Choose
functionals supported on the trivial character of the Cartier
target.  Cartier commutes with etale pullback; the Frobenius twist
permutes characters by squaring and preserves triviality.  Hence
only the trivial block and the inverse pairs
\((\chi,\chi^{-1})\) occur.

Each nontrivial character has odd order and is not two-torsion.
The preceding theorem makes its nine-by-nine cross determinant
a nonzero polynomial on the same nine-dimensional functional
space \(H^0(B^{(2)},M^{(2)})^*\).  The trivial block has a nonempty
rank-eight open: the level-three theta symmetrizer
(Theorem~\ref{thm:ppav-level3-Cartier-symmetric}) identifies its
dual contractions with the alternating Cartier contractions in
Proposition~\ref{prop:DE-A3-ppav-degree-one}, whose generic rank
is eight.  The finitely many nonzero determinant opens and this
rank-eight open intersect, since the affine functional space is
integral over the infinite field \(k\).

At such a functional the total rank is
\(8+18(n-1)/2=9n-1\).  The forms are alternating because
\(C(\omega_A)=0\) on the nonordinary abelian surface and
\(C(s^2\omega_A)=s^{(2)}C(\omega_A)\).
Their space has odd dimension \(9n\), so \(9n-1\) is also the
upper bound.  This proves the generic assertion.
\end{proof}

\begin{corollary}[Every fixed odd principal-cube degree]
\label{cor:fixed-ppav-all-odd-maximal-dual-Cartier}
\status{Proved; independent audits passed.}
For every fixed nonordinary principally polarized abelian surface
\((A,\Theta)\) in characteristic two and every odd integer
\(d\ge1\), the dual Cartier alternating family on
\(H^0(A,\mathcal O_A(3d\Theta))\) has generic rank
\(9d^2-1\).
\end{corollary}
\begin{proof}
Choose a maximal isotropic subgroup \(K\subset A[d]\), of order
\(d^2\), for the theta pairing of \(\mathcal O_A(d\Theta)\).
Maximal isotropic descent gives an etale quotient
\(\pi:A\to B=A/K\) and a principal line bundle \(M_0\) on
\(B\) with \(\pi^*M_0=\mathcal O_A(d\Theta)\).
One may choose \(M_0\) symmetric: if its discrepancy under
\([-1]\) is a character line bundle in \(\ker\pi^*\), adjust
by its unique half in that odd-order kernel.  Pullback is unchanged.
Thus \(\mathcal O_A(3d\Theta)=\pi^*(M_0^3)\).
The surface \(B\) remains nonordinary, since prime-to-two isogenies
preserve the two-rank.  Apply the preceding corollary with
\(n=d^2\).  For \(d=1\), the already cited degree-one contraction
rank is the same assertion.
\end{proof}

\begin{remark}[Scope and independent inputs]
\label{rem:ppav-twisted-Cartier-Moore-scope}
\status{Proved; independent audits passed.}
The last corollary concerns the dual alternating contraction family,
not the full global Lang image.  It resolves the coefficient-vector
identification and outgoing-isogeny direction needed for this rank
statement.  Further local inertia and global generation arguments
are required before claiming new full-monodromy degrees.

The only finite computation used here is the separate, exhaustive
universal Moore classification certificate.  Every new geometric
identity above is proved with relative Frobenius targets, symmetry
compatible theta lifts, and the common parameter fibre explicitly
specified.  The principal f0 bridge is left byte-identical; its
general-character extension is proved here.  The complete independent obstruction/crystal proof and source audit
passed at 01:35:24.622429 UTC on 20 November 2026, and the complete
parent proof and source audit passed at 01:37:50.135808 UTC.
Both checked all seven assertions, the exact theta lifts, every
Frobenius target, the fixed-point argument, and the outgoing
isogeny construction.  No new numerical experiment was needed
for these geometric deductions.  The exact candidate and both
audit records remain retained.
\end{remark}
